\chapter{Conclusão}
\label{sec:conclusion}
O sistema CN2425TF, desenvolvido no âmbito da unidade curricular de Computação na Nuvem, 
demonstra uma aplicação prática e robusta dos serviços oferecidos pela Google Cloud Platform na construção de uma solução escalável, 
tolerante a falhas e orientada a microserviços. A arquitetura proposta permite aos utilizadores submeter imagens para análise de monumentos, 
consultar resultados e obter informação georreferenciada, aproveitando de forma eficaz serviços como Cloud Storage, Firestore, Pub/Sub, 
Cloud Functions, Vision API e Static Maps API.
A utilização do protocolo gRPC como interface entre cliente e servidor revelou-se adequada para garantir comunicações eficientes e estruturadas, 
enquanto o uso de streaming na submissão de imagens permitiu otimizar a transmissão de dados e validar a integridade dos ficheiros submetidos. 
O modelo de dados implementado no Firestore possibilita uma persistência clara, escalável e facilmente consultável dos resultados da análise efetuada.
Adicionalmente, a integração com o Pub/Sub para orquestração assíncrona do processamento das imagens, bem como a separação de responsabilidades entre os diversos componentes, 
contribui para a modularidade e extensibilidade da aplicação. O cliente em linha de comandos permite ao utilizador explorar todas as funcionalidades do sistema de forma simples e eficaz.
